\begingroup
\makeatletter\tagsleft@true\let\veqno\@@leqno\makeatother
\clearpage\SourceLeaf{325b}\ArithmeticHead{318}{L}{ANALYSIS.\quad NACHLASS.}
\NumberedItem{7}
\par\indent Der Zusammenhang zwischen $t$ und $T$, wenn anstatt $T$ S. 176 [der Handschrift, oben S. 312, vergl. die Gl. [5a], S. 314] $T'$ geschrieben
\begin{equation}
T'=T(1+2\lambda\lambda+\lambda^4+2\lambda^6+2\lambda^8+{.}{.}{.}),\tag*{[21]}
\end{equation}
\noindent ist enthalten in
\begin{equation}
\begin{aligned}
\log\frac Tt={}&\frac1{2\pi}\int_0^{2\pi}\left(1-\frac{4\lambda}{1+\lambda\lambda}\cos2e+\frac{4\lambda\lambda}{1+\lambda^4}\cos4e-\frac{4\lambda^3}{1+\lambda^6}\cos6e\,{.}{.}{.}\right)\\
&\qquad.\log(u+\lambda u^{-1}-T).\,de.
\end{aligned}\tag*{[22]}
\end{equation}
\par\indent Der von $T$ unabhängige Theil des Integrals ist
{\small\begin{equation}
\left\{\begin{aligned}
&-\frac{2\lambda\lambda}{1+\lambda\lambda}+\frac12.\frac{2\lambda^4}{1+\lambda^4}-\frac13.\frac{2\lambda^6}{1+\lambda^6}+\text{etc.}\\
&=2\log\left\{\frac{1-\lambda\lambda}{1-\lambda^4}.\frac{1-\lambda^6}{1-\lambda^8}.\frac{1-\lambda^{10}}{1-\lambda^{12}}\,{.}{.}{.}\right\}=2\log\left\{\frac1{1+\lambda\lambda+\lambda^6+\lambda^{12}+\lambda^{20}+\lambda^{30}\,{.}{.}{.}}\right\}.
\end{aligned}\right.\tag*{[23]}
\end{equation}}
\noindent {\small[S. 225]}
\par\indent Setzt man
\begin{equation}
t=AT+BT^3+CT^5+{.}{.}{.},\tag*{[24]}
\end{equation}
\noindent wo nicht zu vergessen ist, dass $A,B,C\,{.}{.}{.}$ und $T$ hier andere Bedeutung haben als oben S. 176 [der Handschrift, oben S. 312], so ist
{\small\begin{equation}
\left\{\begin{aligned}
A&=(1+\lambda\lambda+\lambda^6+\lambda^{12}+\lambda^{20}+\lambda^{30}\,{.}{.}{.})^2\\
&=1+2\lambda\lambda+\lambda^4+2\lambda^6+2\lambda^8+{.}{.}{.},
\end{aligned}\right\}\quad
\begin{gathered}
\text{\textit{Ist die Reihe }}\ast\ \text{\textit{S. 177}}\\
\text{[Gl. [5a], oben S. 314]}
\end{gathered}\tag*{[25]}
\end{equation}}
\noindent wo der Coefficient von $\lambda^{2n}$ einfach mit der Anzahl der Zerlegungen von $4n+1$ in zwei Quadrate zusammenhängt[**)],
\par\bigskip\noindent\rule{.25\linewidth}{.4pt}\par\smallskip
{\small\par\indent [**) Die Handschrift hat »abhängt«.]\par}
\clearpage\SourceLeaf{326}\ArithmeticHead{319}{R}{{\scriptsize ABBILDUNG DER ELLIPSE AUF DEN KREIS.}}
\[\begin{aligned}
B&=-A\left(\frac\lambda{1+\lambda\lambda}+\frac{2\lambda^3}{1+\lambda^4}+\frac{3\lambda^5}{1+\lambda^6}+{.}{.}{.}\right)\\[4pt]
&=-A\left(\frac\lambda{(1-\lambda\lambda)^2}-\frac{\lambda^3}{(1-\lambda^4)^2}+\frac{\lambda^5}{(1-\lambda^6)^2}-{.}{.}{.}\right)\\[4pt]
&=-A\frac{d\log((1+\lambda\lambda)(1+\lambda^4)(1+\lambda^6)\,{.}{.}{.})}{2\,d\lambda}\\[4pt]
&=-A\frac{d\log A}{12\,d\lambda}+\frac{A\,d\log(1-2\lambda\lambda+2\lambda^8-2\lambda^{18}\,{.}{.}{.})}{6\,d\lambda}\\[4pt]
&=-A\{\lambda+\lambda^3+4\lambda^5+\lambda^7+6\lambda^9+4\lambda^{11}+8\lambda^{13}+\lambda^{15}+{.}{.}{.}\},
\end{aligned}\]
\noindent die Coefficienten von $\lambda^{2n+1}$ sind hier aequal dem Aggregat aller ungeraden Divisoren von $2n+2$.
\par\indent Es ist auch
\[\begin{aligned}
B&=-A\left(\frac\lambda{1-\lambda\lambda}+\frac{3\lambda^5}{1-\lambda^6}+\frac{5\lambda^9}{1-\lambda^{10}}+\frac{7\lambda^{13}}{1-\lambda^{14}}\,{.}{.}{.}\right)\\[4pt]
&=+A\frac{d\log(1-\lambda\lambda)(1-\lambda^6)(1-\lambda^{10})\,{.}{.}{.}}{2\,d\lambda},
\end{aligned}\]
\[\frac B{A^3}=\lambda-3\lambda^3+10\lambda^5-29\lambda^7\,{.}{.}{.}\]
\clearpage
\endgroup
